\subsection{Tips \& Tricks}
\shortremark Die \bi{Dichten} können anstelle der Variablen in bspw. der Formel für bedingte W. verwendet werden.

\howto{Randdichten berechnen}{
    \begin{enumerate}
        \item Intervalle zur Integration für Randdichten via Indikatorvariable $\bm{1}_{cond}$ bestimmen. $\bm{1}_x$ verschwindet bei Bestimmung von $f_\cY$ in den Integral-grenzen
        \item $cond$ nach entsprechender Var. auflösen, wenn nötig
        \item Bei Integration Indikatorvariable als konstante betrachten (muss auch integriert werden!)
    \end{enumerate}
}

Die Randdichten sind dann die Dichten von $\cX$, etc.; Für weitere Berechnungen verwendbar, wie bspw. $\E[\cX + \cY]$.

\howto{$cond$ auflösen}{
    In dem Integral sollte jeweils der äussere Integral der sein, für welchen wir volle Information haben.
    Zudem, bspw., für $0 \leq x \leq y^2 \leq 1$, wissen wir zwar für $x$ und $y$ die Grenzen, wählen aber $x$ als inneren Integral,
    weil $\int_0^{y^2}$ sich schöner auflöst
}

\howto{Gemeinsame Dichten finden}{
    Nach Definition muss $f$ die Ableitung von $F$ sein.
    Oder umgekehrt, $\int f \dx \ldots$ sollte $F$ ergeben. Beispiel:
    $T \sim \cU[1, 2], \cX \sim \text{Exp}(\lambda), \cY \sim \text{Exp}(\gamma)$. Dann $f_{T, X, Y}(t, x, y) = \begin{cases}
            \lambda \gamma \exp(-\lambda x - \gamma y) & \text{falls } 1 \leq t \leq 2 \land x, y \geq 0 \\
            0                                          & \text{sonst}
        \end{cases}$
}

\howto{Verteilungen / Dichten finden}{
    Wir suchen $\P[\cZ \leq a]$ (ausser für gem. V, dann s. oben).
    Dann überlegen, wie $\cZ \leq a$ erfüllt werden kann. Beispiel:
    $Z = \min(\cX, \cY)$, dann gesucht: $\P[\cX \leq a \cup \cY \leq a]$ oder $1 - \P[\cZ > a]$.
    Schliesslich $F_Z$ aufstellen mit Bedingungen oder Indikatorvariablen.

    Zudem kann Verteilung der einzelnen Variablen (und v.a. Unabhängigkeit) nützlich werden.

    Falls Dichte, dann Verteilung ableiten.
}

\howto{Multi-Variable Wahrscheinlichkeiten}{
    \bi{Für W. wie} $\P[\cX \leq \cY^3]$ übersetzen wir dies in Integral $\int_B f(x, y) \dx x \dx y$ mit $B = \{(x, y) \divider x \leq y^3 \}$, dann für $y \in [0, 1]$ gilt:
    $\int_0^1 \int_{0}^{y^3} f(x, y) \dx x \dx y$, wobei hier die untere Grenze durch gegebenes $x \geq 0$ entsteht.

    \bi{Für W. im Exponent}, dies beschreibt typ. Repetitionen, also über die Anzahl und W. nachdenken.
}
